%-------- PREAMBLE ----------
\documentclass{SGGW-thesis/SGGW-thesis}
\INZYNIERSKAtrue
\WZIMtrue
%\usepackage{graphicx} % Required for inserting images
\usepackage[utf8]{inputenc} % for encoding
\usepackage[polish]{babel}   % for Polish language support
\usepackage{url}
\usepackage{amsmath}
\usepackage{listings}
\usepackage{xcolor}
\usepackage{hypbmsec}
\usepackage{float}
\usepackage[backend=biber,style=numeric]{biblatex}
\addbibresource{bibliography.bib} % Plik .bib z bibliografią

%-------- PACKAGES SETUP ---------------

\lstset{
    language=Python,                % Określenie języka
    backgroundcolor=\color{white},   % Kolor tła
    basicstyle=\ttfamily\footnotesize, % Czcionka i rozmiar
    numbers=left,                   % Numery wierszy z lewej strony
    numberstyle=\tiny\color{gray},   % Styl numerów
    stepnumber=1,                   % Co ile wierszy numeracja
    numbersep=5pt,                  % Odstęp numerów od kodu
    showspaces=false,                % Nie wyświetlaj spacji
    showstringspaces=false,         % Nie wyświetlaj przestrzeni w ciągach
    showtabs=false,                 % Nie wyświetlaj tabulatorów
    frame=single,                   % Ramka wokół kodu
    tabsize=4,                      % Rozmiar tabulacji
    captionpos=b,                   % Pozycja podpisu (b = na dole)
    breaklines=true,                % Łamanie długich linii
    breakatwhitespace=true,         % Łamanie linii przy spacji
    % escapeinside={\%*}{*)}          % Jeśli chcesz używać LaTeX w kodzie
}

%-------- FIRST PAGE SETUP -----------

\title{System automatycznego sterowania rakietą oparty na symulacjach komputerowych\\i algorytmach sztucznej inteligencji}
\author{Dominik Kubicki 210882}
\date{2025}
\album{Jakub Kindlik 210871}
\thesis{Praca dyplomowa na kierunku}
\course{Informatyka}
\promotor{dra \ inż.\ Pawła Hosera}
\pworkplace{Instytut Informatyki Technicznej\\Katedra Sztucznej Inteligencji}

%----------- DOCUMENT ----------------

\begin{document}
\maketitle
\statementpage

%%----STRESZCZENIE----

\abstractpage
{Temat}
{Streszczenie}
{Słowa kluczowe}
{Temat ang}
{Streszczenie ang}
{Słowa kluczowe ang}

%%----SPIS TREŚCI ----
{
  \doublespacing
  \tableofcontents
}
\startchapterfromoddpage

%----- ROZDZIAŁY ----

\chapter{Wstęp}

\section{Cel pracy}
Celem niniejszej pracy inżynierskiej jest stworzenie systemu sterującego,
opartego na algorytmach samouczących, który będzie w stanie przeprowadzić
stabilny, pionowy lot rakiety.
Aby osiągnąć ten cel pracy, konieczne jest stworzenie środowiska, w którym
algorytm sterujący będzie trenowany.
Stworzenie silnika fizycznego umożliwiającego symulację lotu rakiety oraz
ingerencję w jej lot poprzez zmianę parametrów układu sterującego rakiety jest pośrednim celem niniejszej pracy.

\section{Struktura pracy}

Niniejsza praca inżynierska powstała w wyniku kolaboracji dwóch
współtwórców, dlatego też podzielona jest ona na dwa główne moduły.
\begin{itemize}
    \item Moduł I: Silnik fizyczny symulacji lotów rakietowych. Autor: Jakub
    Kindlik;
    \item Moduł II: Algorytm sterujący rakietą. Autor: Dominik Kubicki;
\end{itemize}
\\
Każdy z modułów jest pełną pracą inżynierską samą w sobie, zawierając
podstawy teoretyczne, przedstawienie wykorzystanych technologii, opis i
dekonstrukcje stworzonego rozwiązania, wnioski, wyniki i opis procesu
powstawania. \\

Łącznie oba projekty-moduły tworzą kompletny system sterujący rakietą,
wytrenowany na silniku fizycznym, co spełnia założony cel pracy.

\\

Pozostałe rozdziały pracy poświęcone są opisowi kwestii dotyczących obu
modułów pracy inżynierskiej, takich jak między innymi wstęp, integracja obu
części projektu, końcowe wnioski wyniki oraz podsumowanie całości rozwiązania
pracy.
\\

Główne moduły pracy są ze sobą tak ściśle związane, że aby zachować spójność,
prace inżynierskie dwóch autorów powstały jako jeden
dokument.
Dzięki temu możliwe jest utrzymanie ciągu logicznego, prowadzącego od
powstania silnika fizycznego przez trenowanie na nim modelu sterującego aż
do integracji obu systemów i wniosków końcowych.

\section{Użyte technologie}
Poszczególne biblioteki i narzędzia wykorzystane podczas pracy nad każdym z
modułów zostały opisane dedykowanych tym modułom rozdziałach pracy.
Jednakże niektóre technologie oraz narzędzia zostały wykorzystane podczas
pracy nad
obydwoma modułami projektu, tworzonymi w ramach niniejszej pracy inżynierskiej.
Są to:

\begin{itemize}
    \item \textbf{Git oraz Github} - narzędzia wykorzystane do kontroli
    wersji projektu.\\ Umożliwiły one wspólną pracę nad kodem projektu oraz
    dostęp do archiwalnych wersji projektu. Git oraz Github zostały wybrane
    ze względu na ich popularność, oraz wsparcie dla wielu języków, a także
    na znajomość tych narzędzi przez autorów pracy.
    \item \textbf{Python 3.13} - język programowania, w którym napisane
    zostały obydwa moduły projektu.\\ Python został wybrany jako język
    programowania w projekcie ze względu na jego prostotę oraz
    znajomość tego języka przez autorów pracy. Ważnym czynnikiem wpływającym
    na wybór tego języka jest również popularność Pythona w dziedzinie
    uczenia maszynowego oraz sztucznej inteligencji. Dzięki wykorzystaniu
    tego języka do obu modułów projektu ich integracja była znacznie prostsza.
    \item \textbf{PyCharm Professional}  - środowisko programistyczne
    wykorzystane do tworzenia kodu projektu.\\PyCharm został wybrany ze
    względu na dedykowane wsparcie dla języka Python, w którym napisany
    został projekt, oraz na znajomość tego narzędzia przez autorów pracy.
\end{itemize}

\\Ponadto, do tworzenia pracy inżynierskiej wykorzystane zostały następujące
technologie umożliwiające stworzenie i edycję niniejszego dokumentu:

\begin{itemize}
    \item \textbf{LaTex} - system składu tekstu wykorzystany do napisania
    niniejszego dokumentu.\\ LaTex został wybrany ze względu na jego
    popularność wśród naukowców oraz inżynierów, prostotę tworzenia pracy
    inżynierskiej w tym narzędziu oraz znajomość tego narzędzia przez autorów pracy.
    \item \textbf{SGGW-thesis} - biblioteka do tworzenia pracy inżynierskiej
    w \latex zgodnie z wymaganiami formatowania pracy inżynierskiej na SGGW.\\ Biblioteka ta została wybrana ze względu na jej prostotę użycia oraz zgodność z wymaganiami formatowania pracy inżynierskiej na SGGW.
\end{itemize}

Zastosowanie tych narzędzi pozwoliło na efektywną pracę nad projektem oraz
na stworzenie tej pracy inżynierskiej. Narzędzia te są szeroko dostpne i
większości darmowe, co pozwoliło na ich wykorzystanie w ramach pracy inżynierskiej.

\section{Przegląd dostępnych rozwiązań. }

Dziedzina hobbystycznej symulacji lotów rakietowych, choć niszowa w kontekście informatyki, inżynierii i fizyki, posiada grono entuzjastów oraz specjalistów, którzy przez lata rozwijali narzędzia i algorytmy umożliwiające symulację takich lotów. Istnieje również duży rynek komercyjnych narzędzi stosowanych przez przedsiębiorstwa oraz wojsko do symulacji lotów rakietowych. Systemy te charakteryzują się wysoką dokładnością, złożonością i są zazwyczaj niedostępne dla hobbystów oraz użytkowników prywatnych. W związku z tym, w niniejszym rozdziale opisane zostały jedynie rozwiązania ogólnodostępne. Poniżej przedstawiono narzędzia i technologie, które umożliwiają symulację lotów rakietowych, a także gotowe silniki fizyczne oraz algorytmy uczenia maszynowego, mogące zostać zastosowane w tym kontekście.

\\Znacznie bardziej popularna dziedzina uczenia maszynowego, oparta na
Reinforcement Learning (RL), oferuje liczne rozwiązania, które mogą zostać wykorzystane do trenowania modeli sztucznej inteligencji odpowiedzialnych za sterowanie lotem rakiety. W tym rozdziale przedstawione zostały najczęściej wykorzystywane z nich — zarówno biblioteki służące do trenowania modeli, jak i narzędzia umożliwiające przeprowadzanie treningów na symulacjach.


\subsection{Silniki symulacji lotów rakietowych.}

Na rynku dostępnych jest wiele modeli silników fizycznych umożliwiających
symulację lotów rakietowych, jak i narzędzi umożliwiających tworzenie
własnych symulacji. Poniżej zaprezentowane zostały najpopularniejsze z nich. \\

Narzędzia służące do ogólnego modelowania symulacji fizycznych:
\begin{itemize}
    \item \textbf{Symulink.} \\Bardzo zaawansowane narzędzie do
    modelowania i prowadzenia symulacji systemów dynamicznych. Symulink
    wchodzi w skład pakietu Matlab. System ten umożliwia elastyczne i
    dokładne modelowanie sił. Narzędzie to bazuje na krokowych modelach
    składania sił, co pozwala na bardzo dokładne modelowanie sił
    działających na obiekt. Symulink jest narzędziem płatnym oraz
    wymagającym dużej wiedzy z zakresu modelowania systemów dynamicznych. \\
    Symulink może zostać wykorzystany do stworzenia bardzo dokładnego modelu
    symulacji lotu rakiety poprzez zamodelowanie sił działających na
    obiekt i symulację ich działania w czasie.


    \item \textbf{ANSYS Fluent.}
    \\Zaawansowane narzędzie do symulacji dynamiki płynów (CFD – Computational Fluid Dynamics) wykorzystywane w inżynierii aerodynamicznej i mechanicznej. ANSYS Fluent umożliwia modelowanie przepływów płynów, wymiany ciepła, oraz oddziaływania sił aerodynamicznych na obiekty w różnych warunkach środowiskowych.
    ANSYS Fluent bazuje na rozwiązaniach równań Naviera-Stokesa, co pozwala na dokładne odwzorowanie zjawisk takich jak opór aerodynamiczny, turbulencje oraz efekty cieplne. Dzięki zaawansowanym algorytmom i funkcjom, Fluent umożliwia analizę złożonych scenariuszy, takich jak przepływ wokół wieloelementowych konstrukcji czy interakcje z gazami spalinowymi.
    Oprogramowanie to jest komercyjne, a jego obsługa wymaga wiedzy z zakresu CFD oraz doświadczenia w konfiguracji symulacji.
    ANSYS Fluent może być wykorzystany do szczegółowego modelowania lotu rakiety poprzez analizę przepływu powietrza wokół korpusu rakiety, co pozwala na optymalizację jej konstrukcji i poprawę stabilności aerodynamicznej.


    \item \textbf{OpenFOAM.}
    \\Otwartoźródłowe oprogramowanie do symulacji CFD, które oferuje zaawansowane możliwości modelowania przepływów płynów, wymiany ciepła oraz innych procesów fizycznych. OpenFOAM opiera się na podejściu siatkowym (mesh-based), co pozwala na precyzyjne odwzorowanie geometrii i rozwiązywanie równań Naviera-Stokesa w ramach analiz aerodynamicznych.
    Główną zaletą OpenFOAM jest jego elastyczność i dostępność – oprogramowanie jest darmowe i można je dostosować do indywidualnych potrzeb poprzez tworzenie własnych modułów i skryptów. Jednak jego obsługa wymaga zaawansowanej wiedzy technicznej oraz umiejętności pracy w środowisku tekstowym, co może być barierą dla początkujących użytkowników. \
    OpenFOAM może zostać wykorzystany do symulacji lotu rakiety poprzez analizę sił aerodynamicznych oraz optymalizację kształtu i konstrukcji rakiety w celu minimalizacji oporu i poprawy wydajności lotu.

\end{itemize}
Dedykowane silniki fizyczne umożliwiające symulację lotów rakietowych:

\begin{itemize}

    \item \textbf{OpenRocket.}
    \\Bezpłatne i otwartoźródłowe narzędzie stworzone specjalnie do symulacji lotów rakiet. OpenRocket jest łatwe w obsłudze, dzięki intuicyjnemu interfejsowi graficznemu, który umożliwia szybkie projektowanie rakiet oraz przeprowadzanie symulacji ich lotu. Program oferuje możliwość modelowania sił aerodynamicznych, takich jak opór powietrza, siła nośna oraz działanie silników rakietowych w różnych warunkach.
    OpenRocket obsługuje różne typy silników, pozwala na analizę stabilności rakiety, trajektorii lotu oraz przewidywanie maksymalnej wysokości i zasięgu. Narzędzie to jest szczególnie popularne wśród hobbystów oraz studentów, dzięki swojej prostocie i dostępności, choć może mieć ograniczenia w przypadku bardziej zaawansowanych symulacji.
    \\Narzędzie to było też w dużym stopniu inspiracją niniejszej pracy
    inżynierskiej.


    \item \textbf{RockSim.}
    \\Komercyjne oprogramowanie do projektowania i symulacji lotów rakiet, dedykowane zarówno dla entuzjastów, jak i profesjonalistów. RockSim umożliwia projektowanie rakiet od podstaw, uwzględniając szczegółowe parametry takie jak geometria, masa komponentów czy rozmieszczenie środka ciężkości i ciśnienia.
    Program oferuje zaawansowane funkcje symulacji aerodynamicznych, takich jak analiza stabilności, przewidywanie trajektorii lotu oraz obliczanie wpływu warunków atmosferycznych (wiatr, temperatura) na zachowanie rakiety. Dzięki swojej wszechstronności, RockSim jest używany zarówno w edukacji, jak i w projektach zaawansowanych konstrukcji rakietowych.


    \item \textbf{RAS Aero.}
    \\RAS Aero to zaawansowane narzędzie do symulacji lotów rakiet dedykowane użytkownikom, którzy potrzebują dużej dokładności w analizie aerodynamicznej. Program oferuje możliwość modelowania trajektorii z uwzględnieniem szczegółowych parametrów aerodynamicznych, takich jak współczynniki oporu czy siły nośnej, oraz wpływu zmiennych warunków atmosferycznych na rakietę.
    RAS Aero jest szczególnie cenione za swoją zdolność do przewidywania zachowania rakiet przy dużych prędkościach oraz podczas przejścia przez bariery dźwiękowe. Oprogramowanie to jest płatne i wymaga większego doświadczenia w pracy z danymi aerodynamicznymi, ale oferuje bardzo dokładne wyniki symulacji.


\end{itemize}

\subsection{Algorytmy uczenia maszynowego oparte o RL.}

W dziedzinie algorytmów uczenia maszynowego opartych o Reinforcement
Learning istnieje wiele narzędzi, jak i gotowych programów, które mogą być
wykorzystane do
uczenia modeli na symulacjach fizycznych.

\\

Najpopularniejsze narzędzia stosowane do uczenia maszyno owego RL w
symulacjach fizycznych to m.in.:


\begin{itemize}
    \item \textbf{TensorFlow.} \\
    TensorFlow to popularna platforma open-source'owa do uczenia maszynowego, która znalazła szerokie zastosowanie w symulacjach fizycznych, w tym w dziedzinie reinforcement learning (RL). Dzięki wsparciu dla zaawansowanych algorytmów optymalizacji, TensorFlow umożliwia trenowanie modeli RL na danych pochodzących z symulacji fizycznych. Framework ten oferuje bogatą funkcjonalność do tworzenia i trenowania sieci neuronowych, a także integrację z innymi narzędziami do analizy danych i wizualizacji. TensorFlow wspiera różnorodne techniki optymalizacji, w tym metody oparte na gradientach, co pozwala na skuteczne rozwiązywanie problemów związanych z dynamiką fizyczną. Jego rozbudowana dokumentacja oraz wsparcie dla rozwoju modeli na dużą skalę czynią go jednym z najczęściej wykorzystywanych narzędzi w branży.

    \item \textbf{PyTorch.} \\
    PyTorch to kolejny popularny framework open-source'owy, który jest szczególnie ceniony wśród badaczy i inżynierów zajmujących się uczeniem maszynowym. Dzięki swojej elastyczności oraz wsparciu dla dynamicznych grafów obliczeniowych, PyTorch doskonale nadaje się do tworzenia i trenowania modeli reinforcement learning w kontekście symulacji fizycznyc. PyTorch jest często wykorzystywany w projektach związanych z robotyką oraz modelowaniem dynamicznych systemów, takich jak rakiety, ponieważ pozwala na łatwą integrację z fizycznymi symulacjami i oferuje szeroki zakres algorytmów RL. PyTorch charakteryzuje się także rozbudowaną społecznością oraz licznymi zasobami edukacyjnymi, co czyni go atrakcyjnym wyborem dla osób rozpoczynających pracę z uczeniem maszynowym.

    \item \textbf{Stable Baselines3.} \\
    Stable Baselines3 to biblioteka open-source'owa zbudowana na bazie PyTorch, która oferuje gotowe implementacje popularnych algorytmów reinforcement learning, takich jak PPO (Proximal Policy Optimization), A2C (Advantage Actor-Critic) czy DQN (Deep Q-Network). Jest to narzędzie szczególnie przydatne w kontekście symulacji fizycznych, gdzie wymagana jest szybka i efektywna implementacja algorytmów RL do testowania i optymalizacji systemów sterowania. Stable Baselines3 jest dobrze udokumentowaną platformą, która ułatwia implementację i eksperymentowanie z różnymi metodami RL. Dzięki swojej prostocie i kompatybilności z popularnymi symulatorami, jak Gazebo czy Unity, biblioteka ta stała się jednym z najczęściej wykorzystywanych narzędzi w projektach związanych z uczeniem maszynowym i symulacjami fizycznymi.
\end{itemize}

\\
Poniżej przedstawione zostały jedne
z popularniejszych oraz częściej wykorzystywanych narzędzi służących do
trenowania modeli sztucznej inteligencji sterujących lotami rakiet.\\

\begin{itemize}
    \item \textbf{Simulink z Aerospace Blockset.} \\
    MATLAB w połączeniu z Simulinkiem oraz dodatkiem Aerospace Blockset to zaawansowane narzędzie umożliwiające symulację lotu rakiet. Oprogramowanie to pozwala na elastyczne modelowanie dynamiki obiektów latających, uwzględniając siły aerodynamiczne, grawitację i inne istotne parametry fizyczne. Aerospace Blockset oferuje gotowe bloki do analizy trajektorii oraz projektowania systemów sterowania. Dodatkowo, MATLAB umożliwia integrację z algorytmami uczenia maszynowego, co pozwala na trenowanie modeli sterowania w oparciu o dane z symulacji. System ten jest płatny i wymaga zaawansowanej wiedzy z zakresu modelowania dynamicznego, jednak oferuje bardzo wysoką precyzję i wsparcie dla projektów naukowych i inżynierskich.

    \item \textbf{Microsoft AirSim.} \\
    Microsoft AirSim to open-source'owy symulator lotu, pierwotnie zaprojektowany dla dronów, ale możliwy do dostosowania do symulacji rakiet. Oprogramowanie to obsługuje realistyczne warunki fizyczne, takie jak zmienne aerodynamiczne i różne środowiska lotu. AirSim integruje się z popularnymi frameworkami uczenia maszynowego, takimi jak TensorFlow czy PyTorch, umożliwiając trenowanie modeli sterowania w oparciu o reinforcement learning. Dzięki elastyczności i wsparciu dla programowania API, użytkownik może tworzyć niestandardowe środowiska do symulacji. AirSim jest darmowy, ale jego zastosowanie wymaga dostosowania środowiska do specyficznych potrzeb użytkownika, co może być czasochłonne.

    \item \textbf{Gazebo + ROS.} \\
    Gazebo to open-source'owy symulator fizyki, który w połączeniu z systemem ROS (Robot Operating System) stanowi potężne narzędzie do modelowania i sterowania obiektami dynamicznymi, takimi jak rakiety. Gazebo oferuje zaawansowaną symulację fizyczną, w tym modelowanie sił aerodynamicznych, kolizji i zmiennych środowiskowych. Integracja z ROS umożliwia tworzenie i implementację algorytmów sterowania opartych na uczeniu maszynowym, takich jak reinforcement learning. Platforma wspiera również analizę danych w czasie rzeczywistym, co czyni ją idealnym narzędziem do testowania i optymalizacji systemów sterowania. Gazebo + ROS jest darmowe, ale wymaga zaawansowanej konfiguracji oraz wiedzy z zakresu programowania i modelowania fizycznego.
\end{itemize}

\subsection{Podsumowanie.}
Jak zostało to przedstawione w tym rozdziale, pomimo dość niszowego
zagadnienia, na rynku dostępnych jest wiele
narzędzi,
które
mogą być
wykorzystane do stworzenia symulacji lotu rakiety oraz trenowania modeli
sztucznej inteligencji. Narzędzia te są w większosci bardzo zaawansowane i
wymagają dużego doświadczenia oraz wiedzy w dziedzinie symulacji, fizyki, jak
i uczenia maszynowego. Często są to także narzędzia płatne. Niewątpliwie
jednak rozwiązania te oferują symulacje o dużej dokładności i
szczegółowości, dzięki zastosowaniu skomplikowanych rozwiązań takich jak CFD
czy analiza trajektorii lotu. Na dokładnych symulacjach trenowane za pomocą
dobrze dobranych algorytmów uczenia maszynowego modele mogą być w stanie
osiągnąć bardzo dobre wyniki w sterowaniu lotem rakiety.\\
\\ W niniejszej pracy inżynierskiej
wykorzystane
zostały własne implementacje silnika fizycznego oraz algorytmu sterującego,
co pozwoliło na pełną kontrolę nad procesem tworzenia i testowania systemu
sterującego rakietą. Efektem tego jest także uproszczony model fizyczny, o
czym mowa jest w kolejnych rozdziałach pracy. Rozwiązanie prezentowane w tej
pracy odbiega od profesjonalnych narzędzi, jednak pozwala na zrozumienie
sposobu działania symulacji lotu rakiety oraz algorytmów sterujących i jest
dobrą bazą pod rozwój bardziej zaawansowanych systemów w przyszłości.
\chapter{Moduł I: Silnik fizyczny symulacji lotów rakietowych.}
\label{ch:silnik-fizyczny-symulacji-lotow-rakietowych.}

Silnik fizyczny do symulacji lotów rakietowych to pierwszy z modułów
niniejszej pracy inżynierskiej.
Za pomocą tego projektu osiągnięty zostaje
cel pośredni pracy - stworzenie środowiska do uczenia algorytmów.
Autorem
tego modułu jest Jakub Kindlik.\\
W tym rozdziale przedstawione zostały teoretyczne podstawy fizyczne, użyte
technologie, oraz
trzy główne komponenty
tego modułu:

\begin{itemize}
    \item Projekt silnika fizycznego.
    \item Model sił działających na rakietę.
    \item Moduł wizualizacyjny symulacji.
\end{itemize}

Omówione zostały także zagadnienia związane ze
sprawdzeniem dokładności symulacji oraz ewolucją projektu.\\

\section{Założenia projektowe.}\label{sec:zalozenia-projektowe.}
W trakcie pracy nad modułem silnika fizycznego symulacji lotów rakietowych poczyniono szereg założeń, które miały na celu ułatwienie implementacji oraz zapewnienie odpowiedniej wydajności symulacji. Poniżej przedstawiono najważniejsze z nich:

\begin{itemize}

\item \textbf{Modelowanie uproszczonych sił działających na rakietę.} \\
W celu zmniejszenia złożoności obliczeniowej model sił działających na rakietę zostanie odpowiednio uproszczony. Uwzględnione zostaną jedynie kluczowe siły, takie jak grawitacja, opór powietrza i siła ciągu silnika. Takie podejście umożliwia łatwiejsze zamodelowanie systemu i redukuje zapotrzebowanie na zasoby obliczeniowe.

\item \textbf{Optymalizacja pod kątem pracy na komputerze osobistym.} \\
Symulacje będą wykonywane na komputerze osobistym, dlatego silnik fizyczny
musi być zoptymalizowany pod kątem wydajności, jak i okrojony w kwestii
skomplikowanych obliczeń. Ograniczenie
zasobożerności
jest konieczne, aby zapewnić płynne działanie nawet na sprzęcie o przeciętnych parametrach.

\item \textbf{Uproszczona geometria rakiety.} \\
W celu ograniczenia złożoności modelu fizycznego rakieta będzie miała kształt walca i zostanie pozbawiona skrzydeł. Taka decyzja pozwala na skupienie się na kluczowych aspektach dynamiki lotu bez nadmiernego komplikowania obliczeń.

\item \textbf{Sterowanie poprzez obracany silnik.} \\
Układ sterujący rakietą będzie się opierał na możliwości obracania silnika,
co bezpośrednio wpłynie na trajektorię lotu. Taki model jest prosty do
zamodelowania, a jednocześnie pozwala na sterowanie rakietą przez algorytmy
samouczące, poprzez zmianę jednego parametru.

\item \textbf{Możliwość modyfikacji parametrów symulacji w czasie rzeczywistym.} \\
Symulacja zostanie zaprojektowana w sposób umożliwiający ingerencję w jej
parametry w trakcie trwania. W szczególności modyfikowanym parametrem
powinien być kąt nachylenia silnika. Taka funkcjonalność pozwala na
uczenie algorytmów sterujących w czasie rzeczywistym.

\item \textbf{Możliwość wizualizacji symulacji.} \\
Symulacja będzie wyposażona w opcjonalną funkcję wizualizacji, co umożliwi analizę trajektorii i zachowania rakiety w sposób intuicyjny. Jednak wizualizacja nie będzie wymagana, co pozwoli ograniczyć zużycie zasobów obliczeniowych w trakcie symulacji.

\end{itemize}

Założenia te konieczne są do stworzenia silnika fizycznego, który będzie w
stanie współpracować z modułem algorytmu sterującego rakietą, jednocześnie
będąc na tyle prostym, aby nie obciążać zasobów komputera osobistego.


\section{Teoretyczne podstawy fizyczne.}\label{sec:teoretyczne-podstawy.}

Tworzenie silnika symulacji lotów rakietowych wymaga zrozumienia podstawowych
zasad fizyki odnoszących się do ruchu rakiety w atmosferze.


\section{Użyte technologie.}\label{sec:uzyte-technologie.}

Do stworzenia silnika fizycznego symulacji lotów rakietowych wykorzystane zostały następujące technologie:

\begin{itemize}
    \item \textbf{Python 3.13} - język programowania
    \item zestaw bibliotek
    \begin{itemize}
        \item bibioteka1....
    \end{itemize}
\end{itemize}

\section{Projekt silnika fizycznego.}\label{sec:projekt-silnika-fizycznego.}
\section{Model sił działających na rakietę.}\label{sec:model-si-dziaajacych-na-rakiete.}
\section{Moduł wizualizacyjny symulacji.}\label{sec:modul-wizualizacyjny-symulacji.}
\section{Dokładność symulacji.}\label{sec:dokadnosc-symulacji}
\section{Ewolucja projektu i wnioski.}\label{sec:ewolucja-projektu-i-wnioski}
\chapter{Moduł II: Algorytm sterujący rakietą.}\

Algorytm sterujący rakietą jest drugim modułem projektu.
Jego celem jest stworzenie algorytmu, który będzie w stanie samodzielnie sterować rakietą w
symulacji. Zadaniem algorytmu będzie pionizacja trajektorii lotu rakiety.
Algorytm będzie uczony na silniku symulacyjnym będącym efektem pracy nad
modułem I.\\
Autorem tego modułu jest Dominik Kubicki.\\
Niniejszy rozdział opisuje założenia projektowe, teoretyczne podstawy, użyte
technologie oraz wybór odpowiedniego algorytmu uczenia maszynowego, a także
sposoby i efekty uczenia algorytmu.\\


\section{Założenia i cele projektowe.}\label{sec:założenia-projektowe-2.}

W trakcie pracy nad algorytmem sterującym rakietą poczynione zostały następujące założenia:

\begin{itemize}
    \item Algorytm sterujący rakietą będzie miał za zadanie pionizację
    trajektorii lotu rakiety, czyli doprowadzenie rakiety do pionu w rożnych
    warunkach atmosferycznych. Bardziej skomplikowane trajektorie lotu nie
    są obiektem badań w tej pracy.
    \item Algorytm będzie uczony na silniku symulacyjnym będącym efektem pracy nad modułem I.
    \item Algorytm będzie wykorzystywał algorytmy sztucznej inteligencji do
    nauki, która będzie przebiegała w sposób dynamiczny w trakcie lotu rakiety.
    \item Konieczne jest stworzenie komponentu umożliwiającego komunikację
    pomiędzy silnikiem symulacyjnym a algorytmem sterującym rakietą.
\end{itemize}

Założenia te pozwalają na usystematyzowanie pracy nad algorytmem oraz
ubranie go w ramy teoretyczne, upraszczając skomplikowanie problemu.\\


\section{Teoretyczne podstawy uczenia modeli sztucznej inteligencji.}
\label{ch:teoretyczne-podstawy-uczenia-modeli-sztucznej-inteligencji.}


W niniejszym rozdziale przedstawione zostały kompleksowe podstawy teoretyczne
dotyczące uczenia modeli sztucznej inteligencji, ze szczególnym
uwzględnieniem metod stosowanych w sterowaniu lotami rakietowymi. Omówione
zostały różne paradygmaty uczenia maszynowego, ich matematyczne fundamenty oraz praktyczne aspekty implementacyjne w kontekście systemów autonomicznego sterowania. Rozdział stanowi teoretyczną podbudowę dla przedstawionych w dalszej części pracy rozwiązań praktycznych.


\subsection{Podejścia do uczenia maszynowego.}
\label{sec:rodzaje-uczenia-maszynowego.}

Współczesne systemy sztucznej inteligencji wykorzystują różne paradygmaty
uczenia maszynowego, z których każdy znajduje zastosowanie w innych
scenariuszach sterowania rakietą. Poniżej przedstawiono charakterystykę i
obszary zastosowań trzech przewodnich podejść do uczenia maszynowego.\\

\begin{itemize}
\item \textbf{Uczenie nadzorowane (Supervised Learning)} \\
W uczeniu nadzorowanym model jest trenowany na podstawie zestawu danych wejściowych i odpowiadających im poprawnych wyników (etykiet). W przypadku sterowania rakietą, dane wejściowe mogą obejmować parametry lotu (np. prędkość, wysokość, kąt nachylenia), a etykiety to optymalne działania sterujące (np. pożądany kąt nachylenia silnika). Algorytm uczy się mapować dane wejściowe na odpowiednie działania, minimalizując błąd pomiędzy przewidywaniami a rzeczywistymi etykietami. Przykłady modeli stosowanych w uczeniu nadzorowanym to sieci neuronowe, drzewa decyzyjne czy modele regresji.

\item \textbf{Uczenie przez wzmacnianie (Reinforcement Learning)} \\
W podejściu tym model uczy się poprzez interakcję ze środowiskiem (w tym przypadku symulatorem lotu), otrzymując w każdym kroku sygnał nagrody za dobre działania. W prezentowanym rozwiązaniu, w każdym kroku symulacji model otrzymuje aktualny stan rakiety i zwraca decyzję o kącie wychylenia silnika, a następnie otrzymuje ocenę tej decyzji w postaci funkcji nagrody. Takie podejście pozwala na uczenie się optymalnej strategii sterowania bez potrzeby posiadania gotowych przykładów "poprawnych" decyzji.

\item \textbf{Uczenie nienadzorowane (Unsupervised Learning)} \\
W uczeniu nienadzorowanym model analizuje dane bez dostępu do etykiet, skupiając się na odkrywaniu ukrytych wzorców lub struktur w danych. Choć mniej bezpośrednio przydatne w sterowaniu rakietą, metody te mogą być wykorzystane np. do grupowania podobnych stanów lotu lub redukcji wymiarowości danych. Przykłady algorytmów to k-means czy autoenkodery.
\end{itemize}


W niniejszej pracy zastosowano podejście z zakresu uczenia przez wzmacnianie, które szczególnie dobrze sprawdza się w problemach sterowania dynamicznego, gdzie pełna znajomość modelu fizycznego jest trudna do uzyskania, a tradycyjne metody kontrolne mogą być niewystarczająco elastyczne.

\subsection{Algorytmy uczenia przez wzmacnianie}
W kontekście sterowania rakietą szczególnie istotne są następujące klasy algorytmów uczenia przez wzmacnianie, każda oferująca unikalne podejście do problemu optymalizacji strategii sterowania:

\paragraph{Metody gradientu polityki (Policy Gradient Methods)}\\
Stanowią podstawową rodzinę algorytmów bezpośrednio optymalizujących funkcję nagrody poprzez aktualizację parametrów polityki. Ich matematyczną podstawę opisuje równanie:

\begin{equation}
    \nabla_\theta J(\theta) = \mathbb{E}_\pi\left[\nabla_\theta \log \pi_\theta(a|s) Q^\pi(s,a)\right]
    \label{eq:policy_gradient}
\end{equation}

gdzie:
\begin{itemize}
    \item $\theta$ reprezentuje parametry polityki sterowania
    \item $\pi_\theta(a|s)$ to prawdopodobieństwo podjęcia akcji $a$ w stanie $s$
    \item $Q^\pi(s,a)$ oznacza oczekiwaną nagrodę dla danej pary stan-akcja
\end{itemize}

W praktyce rakietowej, metody te pozwalają na bezpośrednie uczenie się
strategii sterowania bez konieczności modelowania pełnej funkcji wartości.
Ich główną zaletą jest zdolność do obsługi ciągłych przestrzeni akcji, co jest kluczowe dla precyzyjnego sterowania rakietą. Wyzwaniem pozostaje jednak wysoka wariancja estymatorów gradientu, co może prowadzić do niestabilności procesu uczenia.\\

\paragraph{Architektura Aktor-Krytyk (Actor-Critic)}\\
Stanowi hybrydowe podejście łączące zalety metod opartych na funkcji wartości i metod gradientu polityki. Składa się z dwóch współpracujących komponentów:

\begin{itemize}
    \item \textbf{Aktor (policy network)} - odpowiedzialny za generowanie akcji sterujących na podstawie aktualnego stanu rakiety. Działa jako "pilot", podejmujący decyzje w czasie rzeczywistym.

    \item \textbf{Krytyk (value network)} - pełni rolę "instruktora", oceniającego jakość podjętych decyzji poprzez estymację oczekiwanej nagrody. Jego wyjście służy do redukcji wariancji aktualizacji polityki.
\end{itemize}

W systemach sterowania rakietą, architektura ta szczególnie dobrze sprawdza się dzięki równowadze między eksploracją nowych strategii (aktor) a efektywnym wykorzystaniem zebranych doświadczeń (krytyk). Implementacje takie jak A2C (Advantage Actor-Critic) czy A3C (Asynchronous Advantage Actor-Critic) dodatkowo poprawiają efektywność poprzez wykorzystanie równoległych środowisk.

\paragraph{Projektowanie funkcji nagrody}\\
Kluczowym aspektem skutecznego zastosowania RL w sterowaniu rakietą jest odpowiednie sformułowanie funkcji nagrody, która musi precyzyjnie odzwierciedlać cele misji. Typowa funkcja nagrody dla problemu stabilizacji rakiety może uwzględniać następujące komponenty:

\begin{itemize}
    \item \textbf{Śledzenie trajektorii}: Nagroda za zmniejszanie odległości od zadanej ścieżki lotu, często wyrażana jako ujemna norma różnicy między aktualną a docelową pozycją ($\|p_t-p_{target}\|$)

    \item \textbf{Stabilność orientacji}: Kara za odchylenia od pożądanej orientacji, mierzona poprzez kąty Eulera (roll, pitch, yaw)

    \item \textbf{Efektywność energetyczna}: Nagroda za oszczędne zużycie paliwa, często proporcjonalna do kwadratu zastosowanego sterowania ($a_t^2$)

    \item \textbf{Bezpieczeństwo}: Ostre kary za przekroczenie dopuszczalnych parametrów lotu (np. zbyt duże przyspieszenia, kąty natarcia)
\end{itemize}

Przykładowa funkcja nagrody łącząca te elementy:

\begin{equation}
    r_t = -w_1\|p_t-p_{target}\| - w_2\|\theta_t\| - w_3 a_t^2 - w_4 \mathbb{I}_{\text{awaria}}
    \label{eq:reward_function}
\end{equation}

gdzie wagi $w_i$ determinują względne znaczenie każdego komponentu, a $\mathbb{I}_{\text{awaria}}$ to funkcja indykatora sytuacji awaryjnej. Odpowiednie zbalansowanie tych składowych jest kluczowe - zbyt mocny nacisk na stabilizację może prowadzić do zachowań zbyt konserwatywnych, podczas gdy nadmierna nagroda za śledzenie trajektorii może powodować niebezpieczne oscylacje.

\paragraph{Wyzwania w zastosowaniu do sterowania rakietą}
Implementacja systemów RL dla sterowania rakietą wiąże się z szeregiem istotnych wyzwań technicznych:

\begin{itemize}
    \item \textbf{Wysoka wymiarowość przestrzeni stanów}: Pełny stan dynamiczny rakiety może obejmować kilkanaście zmiennych (pozycja, prędkość, orientacja, prędkość kątowa, parametry silnika), co wymaga zaawansowanych technik aproksymacji funkcji wartości.

    \item \textbf{Niestacjonarność dynamiki}: Zmieniająca się masa rakiety (w wyniku spalania paliwa) oraz zmienne warunki atmosferyczne prowadzą do ciągłych zmian parametrów systemu, wymagając algorytmów zdolnych do adaptacji.

    \item \textbf{Ryzyko katastrofalnych awarii}: Błędy w fazie uczenia mogą prowadzić do niestabilności i utraty kontroli, co w rzeczywistych systemach jest niedopuszczalne. Konieczne są mechanizmy bezpieczeństwa i symulacje w pętli sprzężenia zwrotnego (hardware-in-the-loop).

    \item \textbf{Wymagania wierności symulacji}: Niedokładności w modelu fizycznym mogą prowadzić do uczenia się strategii nieprzenoszących się na rzeczywisty system. Rozwiązaniem jest wykorzystanie technik randomizacji domeny (domain randomization) lub uczenia meta-RL.
\end{itemize}

Pomimo tych wyzwań, ostatnie postępy w dziedzinie głębokiego RL pokazują, że odpowiednio zaprojektowane systemy są w stanie osiągać nadludzką precyzję w sterowaniu złożonymi systemami dynamicznymi, włączając w to problematykę lotów rakietowych. Kluczem jest staranne połączenie solidnych podstaw teoretycznych z praktycznymi technikami stabilizacji uczenia i walidacji bezpieczeństwa.

\subsection{Modele stosowane w sterowaniu rakietą.}
\label{sec:modele-stosowane-w-sterowaniu-rakieta.}

Wybór odpowiedniej architektury modelu ma kluczowe znaczenie dla skuteczności systemu sterowania. Poniżej przedstawiono analizę najważniejszych podejść.

\subsubsection{Modele dla uczenia nadzorowanego}
\begin{itemize}
\item \textbf{Sieci neuronowe MLP (Multi-Layer Perceptron)} \
Architektura MLP sprawdza się w zadaniach regresji parametrów sterowania. Wielowarstwowa struktura z nieliniowymi funkcjami aktywacji pozwala na modelowanie złożonych zależności między stanem rakiety a optymalnymi komendami sterującymi. Typowa implementacja zawiera 3-5 warstw ukrytych o rozmiarze 64-256 neuronów.

\item \textbf{Regresja Gaussa (Gaussian Process Regression)} \
Metoda bayesowska szczególnie przydatna gdy dostępna jest ograniczona liczba danych treningowych. Pozwala na estymację niepewności predykcji, co może być cenne w krytycznych fazach lotu. Wadą jest wyższa złożoność obliczeniowa.
\end{itemize}

\subsubsection{Modele dla uczenia nienadzorowanego}
\begin{itemize}
\item \textbf{Autoenkodery wariacyjne (VAE)} \
Pozwalają na kompresję wysokowymiarowych danych sensorycznych do reprezentacji latentnej, zachowując kluczowe informacje o stanie rakiety. Mogą służyć jako wstępny etap przetwarzania przed właściwym algorytmem sterującym.

\item \textbf{GMM (Gaussian Mixture Models)} \
Użyteczne do klasteryzacji stanów rakiety i identyfikacji typowych konfiguracji. Mogą wspomagać system sterowania poprzez szybką klasyfikację sytuacji awaryjnych.
\end{itemize}

\subsubsection{Modele dla uczenia przez wzmacnianie}
\begin{itemize}
\item \textbf{DDPG (Deep Deterministic Policy Gradient)} \
Algorytm łączący podejście policy gradient z Q-learningiem, szczególnie skuteczny dla ciągłych przestrzeni akcji. Wymaga jednak starannego dostrojenia hiperparametrów i może być niestabilny.

\item \textbf{SAC (Soft Actor-Critic)} \
Nowoczesny algorytm maksymalizujący zarówno nagrodę jak i entropię polityki, co prowadzi do bardziej eksploracyjnych zachowań. Dobrze radzi sobie z zadaniami o ciągłej przestrzeni akcji.
\end{itemize}


\subsection{Wybrany model - PPO (Proximal Policy Optimization)}
Wśród różnych algorytmów uczenia przez wzmacnianie, Proximal Policy Optimization (PPO) wyróżnia się jako szczególnie odpowiedni wybór dla problemów sterowania rakietą. Opracowany w 2017 roku przez zespół badawczy OpenAI pod kierownictwem Johna Schulmana, PPO szybko stał się złotym standardem w zastosowaniach wymagających stabilnego uczenia w ciągłych przestrzeniach akcji - dokładnie taki przypadek mamy w sterowaniu rakietą.

\paragraph{Charakterystyka teoretyczna}\\
Podstawą teoretyczną PPO jest zmodyfikowana funkcja kosztu, która wprowadza inteligentne ograniczenie aktualizacji polityki. Matematycznie wyraża się to wzorem:

\begin{equation}
L^{CLIP}(\theta) = \mathbb{E}_t[\min(r_t(\theta)\hat{A}_t, \text{clip}(r_t(\theta), 1-\epsilon, 1+\epsilon)\hat{A}_t)]
\end{equation}

gdzie poszczególne składniki mają kluczowe znaczenie:
\begin{itemize}
    \item $r_t(\theta) = \frac{\pi_\theta(a_t|s_t)}{\pi_{\theta_{old}}(a_t|s_t)}$ - stosunek prawdopodobieństw wskazujący, jak bardzo nowa polityka różni się od starej
    \item $\hat{A}_t$ - estymata funkcji przewagi, mierząca względną jakość podjętej akcji
    \item $\epsilon$ - parametr ograniczający (zwykle 0.1-0.3), który determinuje maksymalny dopuszczalny krok aktualizacji
\end{itemize}

Mechanizm ten zapewnia, że aktualizacje polityki są wystarczająco duże, by efektywnie się uczyć, ale na tyle małe, by zachować stabilność procesu treningowego - kluczowa właściwość w bezpieczeństwie krytycznych aplikacjach rakietowych.

\paragraph{Zalety w kontekście sterowania rakietą}\\
PPO oferuje szereg unikalnych korzyści w zastosowaniach aerokosmicznych:

\begin{itemize}
    \item \textbf{Bezpieczeństwo i stabilność uczenia}: Mechanizm clip
    skutecznie zapobiega zbyt radykalnym zmianom polityki, które w przypadku
    rakiet mogłyby prowadzić do katastrofalnych awarii. \
    \item \textbf{Efektywność obliczeniowa}: W przeciwieństwie do swojego poprzednika TRPO, PPO nie wymaga kosztownych obliczeń macierzy Hessego, co pozwala na szybsze iteracje treningowe. W praktyce oznacza to możliwość przeprowadzenia pełnego treningu w rozsądnym czasie nawet na pojedynczym GPU.

    \item \textbf{Elastyczność architektoniczna}: PPO doskonale radzi sobie
    zarówno z prostymi układami sterowania (np. stabilizacja jednej osi),
    jak i złożonymi manewrami 3D. \

    \item \textbf{Scalalność}: Algorytm łatwo zrównoleglić na wiele
    instancji środowiska, co znacząco przyspiesza zbieranie doświadczeń.\
\end{itemize}

\paragraph{Ograniczenia i wyzwania}\\
Mimo licznych zalet, PPO nie jest pozbawiony wad, które należy uwzględnić w projektowaniu systemu:

\begin{itemize}
    \item \textbf{Eksploracja przestrzeni stanów}: W porównaniu do algorytmów jak SAC (Soft Actor-Critic), PPO może mieć trudności z wyjściem z lokalnych minimów w szczególnie złożonych scenariuszach. W praktyce rakietowej często wymaga to wprowadzenia dodatkowych mechanizmów eksploracji.

    \item \textbf{Wrażliwość na funkcję nagrody}: Podobnie jak wszystkie metody RL, PPO jest bardzo czuły na projekt funkcji nagrody. Niewłaściwie zbalansowane wagi poszczególnych składowych mogą prowadzić do nieoczekiwanych strategii sterowania.

    \item \textbf{Precyzja kontroli}: Choć PPO generalizuje lepiej niż wiele innych metod, w niektórych zastosowaniach może ustępować specjalistycznym kontrolerom (np. PID z adaptacyjnymi parametrami) w zadaniach wymagających ekstremalnej precyzji.
\end{itemize}


Podsumowując, PPO stanowi doskonały kompromis między stabilnością, efektywnością i praktyczną użytecznością w problemach sterowania rakietą. Jego unikalne właściwości, szczególnie mechanizm ograniczonej aktualizacji polityki, czynią go preferowanym wyborem dla zastosowań wymagających zarówno bezpieczeństwa, jak i adaptacyjności. Mimo pewnych ograniczeń, odpowiednio zaimplementowany i dostrojony algorytm PPO jest w stanie osiągać wyniki porównywalne z zaawansowanymi kontrolerami specjalistycznymi, oferując przy tym znaczną elastyczność.


\subsection{Podsumowanie}

Przedstawione w rozdziale teoretyczne podstawy uczenia maszynowego stanowią fundament dla opracowanego systemu sterowania rakietą. Wybór metody uczenia przez wzmacnianie, a w szczególności algorytmu PPO, podyktowany został jego unikalnymi właściwościami w kontekście problemów sterowania dynamicznego. Zaprezentowane modele i algorytmy, poparte solidnymi podstawami matematycznymi, tworzą kompleksową teoretyczną podbudowę dla praktycznych rozwiązań przedstawionych w kolejnych rozdziałach pracy.

W szczególności, analiza zalet i ograniczeń poszczególnych podejść pozwoliła na świadomy wybór architektury systemu, który łączy w sobie stabilność uczenia, efektywność obliczeniową i praktyczną przydatność w rzeczywistych zastosowaniach aerokosmicznych.


\section{Użyte technologie.}\label{sec:użyte-technologie-2.}
Do stworzenia algorytmu sterującego rakietą wykorzystane zostały następujące technologie:

\begin{itemize}
    \item 1
    \item 2
\end{itemize}


\section{Proces uczenia i optymalizacja.}
\section{Ewolucja projektu i wnioski.}
\chapter{Integracja komponentów: Silnik fizyczny i sztuczna inteligencja}
\section{Interakcja między silnikiem fizycznym a modelem AI}
\section{Synchronizacja działania obu modułów}
\section{Rozwój interfejsu użytkownika}
\chapter{Wyniki działania i wnioski}
\section{Ewaluacja działania silnika fizycznego}
\section{Ocena skuteczności modelu sztucznej inteligencji}
\section{Analiza osiągniętych rezultatów i ich implikacje}
\chapter{Podsumowanie}
\section{Wnioski końcowe}
\section{Perspektywy rozwoju i dalsze badania}


%---- BIBLIOGRAFIA ----
\printbibliography

\beforelastpage

\end{document}
